
% !TeX root=../main.tex
\chapter{مروری بر مطالعات انجام شده}
%\thispagestyle{empty} 
\section{مقدمه}

ناوبری خودران ربات‌های متحرک چرخ‌دار، به‌ویژه در محیط‌های دینامیکی، یکی از مسائل بنیادی و چالش‌برانگیز رباتیک است. گسترش ربات‌های خدمت‌رسان، وسایل نقلیه سبک خودران و سامانه‌های حمل‌ونقل هوشمند، نیاز به حرکت ایمن، روان و قابل‌اعتماد بدون دخالت مستقیم انسان را افزایش داده است؛ درحالی‌که تغییرات مداوم محیط و رفتارهای غیرقابل پیش‌بینی، تصمیم‌گیری ناوبری را دشوار می‌کنند%
\cite{Zhu2021}.

در محیط دینامیکی، موانع ثابت تنها بخشی از چالش‌اند و پیچیدگی اصلی از موانع متحرک با سرعت و مسیر متغیر، عامل‌های تصمیم‌گیر مستقل مانند انسان‌ها و سایر ربات‌ها و تغییرات موقتی ساختار محیط ناشی می‌شود. ربات باید علاوه بر اجتناب از برخورد، پیامد زمانی تصمیم خود را نیز در نظر بگیرد و پیش‌بینی کند که آیا مسیر انتخاب‌شده با جابه‌جایی آینده موانع همچنان ایمن خواهد بود. ازاین‌رو، ناوبری از مسئله‌ای صرفاً هندسی به مسئله‌ای پویا و وابسته به زمان تبدیل می‌شود%
\cite{s25113394}.

روش‌های کلاسیک که بر مدل‌های صریح، قوانین از پیش تعریف‌شده و توابع هزینه ثابت متکی‌اند، اغلب محیط را ایستا یا تغییرات آن را ساده و قابل پیش‌بینی فرض می‌کنند. در حضور رفتارهای غیرقابل پیش‌بینی و تعاملات چندعاملی، نیاز به تنظیم دستی افزایش می‌یابد و عملکرد الگوریتم به شرایط خاص محیط وابسته می‌شود؛ در نتیجه، تعمیم‌پذیری این روش‌ها در محیط‌های واقعی و متغیر محدود است%
\cite{s23249672}.

ویژگی‌های فیزیکی پلتفرم نیز پیچیدگی مسئله را افزایش می‌دهند. اسکوتر چهار چرخ محرک در مقایسه با ربات‌های دیفرانسیلی سبک آزمایشگاهی، جرم و اینرسی بیشتر، فاصله ترمز طولانی‌تر و محدودیت‌های حرکتی متفاوتی دارد. بنابراین، خطاهای کوچک در انتخاب مسیر یا سرعت می‌توانند پیامدهای جدی‌تری برای ایمنی و پایداری حرکت داشته باشند و بسیاری از فرضیات ساده‌کننده ربات‌های کوچک برای این پلتفرم معتبر نیستند.

پیشرفت «یادگیری تقویتی عمیق%
\LTRfootnote{Deep Reinforcement Learning -- DRL}»
امکان یادگیری سیاست ناوبری از طریق تعامل مستقیم با محیط و بدون مدل‌سازی صریح همه حالات را فراهم کرده است. مرورهای اخیر نشان می‌دهند که DRL به یکی از جریان‌های اصلی ناوبری ربات‌های متحرک در محیط‌های پیچیده و دینامیکی تبدیل شده است%
\cite{Zhu2021,s25113394}.
بااین‌حال، بسیاری از مطالعات موجود بر پلتفرم‌های ساده، محیط‌های محدود یا سناریوهای کم‌پیچیدگی متمرکزند و ناوبری مبتنی بر DRL برای پلتفرم‌های واقعی‌تر با قیود دینامیکی جدی، به‌ویژه در محیط‌های شلوغ، کمتر بررسی شده است.

بر این اساس، هدف این پایان‌نامه طراحی و ارزیابی سامانه‌ای مبتنی بر یادگیری تقویتی عمیق برای ناوبری خودران اسکوتر چهار چرخ محرک است تا با لحاظ دینامیک پلتفرم و تغییرات محیط، فرمان‌هایی ایمن، هموار و کارا در حضور موانع متحرک و شرایط غیرقابل پیش‌بینی تولید کند.

\section{روش‌های کلاسیک ناوبری ربات‌های متحرک}

پیش از گسترش رویکردهای مبتنی بر یادگیری، مسئله ناوبری ربات‌های متحرک عمدتاً با استفاده از روش‌های کلاسیک حل می‌شد. این روش‌ها معمولاً بر پایه مدل‌های صریح محیط، قوانین هندسی، جست‌وجو در فضای حالت، توابع هزینه یا حل مسائل بهینه‌سازی طراحی می‌شوند. مزیت اصلی آن‌ها شفافیت، قابلیت تحلیل، امکان پیاده‌سازی بلادرنگ و رفتار نسبتاً قابل پیش‌بینی است. به همین دلیل، بسیاری از سامانه‌های عملی ناوبری هنوز نیز از ترکیبی از این روش‌ها، به‌ویژه برای برنامه‌ریزی مسیر کلی و اجتناب محلی از برخورد، استفاده می‌کنند.

با این حال، روش‌های کلاسیک معمولاً بر فرض‌هایی مانند دسترسی به نقشه نسبتاً دقیق، قابلیت پیش‌بینی محیط، مدل دینامیکی مشخص و تنظیم مناسب پارامترها متکی هستند. این فرض‌ها در محیط‌های دینامیکی، شلوغ و ناقصاً مشاهده‌پذیر همواره برقرار نیستند. از این رو، بررسی این روش‌ها برای روشن شدن محدودیت‌های آن‌ها و توجیه حرکت به سمت رویکردهای یادگیری‌محور ضروری است.

\subsection{برنامه‌ریزی مسیر سراسری مبتنی بر گراف و جست‌وجو}

یکی از قدیمی‌ترین و پرکاربردترین خانواده‌های روش‌های کلاسیک ناوبری، روش‌های مبتنی بر گراف و جست‌وجو هستند. در این رویکردها، محیط معمولاً به‌صورت یک شبکه، گراف یا نقشه اشغال نمایش داده می‌شود و مسئله ناوبری به یافتن مسیر کم‌هزینه میان موقعیت اولیه و هدف تبدیل می‌گردد. الگوریتم‌هایی مانند \lr{Dijkstra}، \lr{A*}، \lr{D*} و \lr{D* Lite} از شناخته‌شده‌ترین روش‌های این دسته هستند%
\cite{Dijkstra1959,Hart1968AStar,koenig2005dstar}.
در این روش‌ها، هزینه مسیر می‌تواند بر اساس طول مسیر، فاصله از موانع، زمان حرکت یا ترکیبی از معیارهای هندسی تعریف شود.

روش‌های جست‌وجومحور به‌ویژه برای محیط‌های ایستا و دارای نقشه مشخص مناسب هستند، زیرا مسئله ناوبری را به جست‌وجوی مسیر کم‌هزینه در یک نمایش گرافی از محیط تبدیل می‌کنند. برای مثال، الگوریتم \lr{A*} با استفاده از تابع ابتکاری می‌تواند فرآیند جست‌وجو را نسبت به روش‌های جست‌وجوی کور هدفمندتر کند. همچنین روش‌هایی مانند \lr{D*} و \lr{D* Lite} برای بازبرنامه‌ریزی مسیر در محیط‌هایی که بخشی از نقشه به‌تدریج آشکار می‌شود، توسعه یافته‌اند و در ربات‌های متحرک کاربرد زیادی دارند%
\cite{koenig2005dstar,ferguson2008motion_planning}.

با وجود این مزایا، محدودیت اصلی این دسته از روش‌ها آن است که معمولاً مسیر را در سطح هندسی و سراسری تولید می‌کنند و به‌تنهایی تضمین نمی‌کنند که مسیر تولیدشده با دینامیک واقعی پلتفرم، محدودیت شتاب، محدودیت ترمز و قیود ایمنی لحظه‌ای سازگار باشد. علاوه بر این، در محیط‌های دینامیکی که موانع متحرک به‌صورت پیوسته موقعیت خود را تغییر می‌دهند، مسیر سراسری ممکن است به‌سرعت نامعتبر شود و سیستم به بازبرنامه‌ریزی مکرر یا ترکیب با یک برنامه‌ریز محلی نیاز پیدا کند%
\cite{LaValle2006Planning,choset2005principles}.
این مسئله برای پلتفرم‌هایی مانند اسکوتر چهار چرخ محرک که دارای اینرسی و فاصله توقف قابل توجه هستند، اهمیت بیشتری دارد.

\subsection{روش‌های نمونه‌برداری در فضای پیکربندی}

دسته دیگری از روش‌های کلاسیک ناوبری، روش‌های نمونه‌برداری در فضای پیکربندی هستند. در این رویکردها، به‌جای جست‌وجوی کامل یا گسسته‌سازی منظم فضای پیکربندی، مجموعه‌ای از پیکربندی‌های مجاز نمونه‌برداری شده و سپس به‌صورت گراف یا درخت به یکدیگر متصل می‌شوند. روش‌هایی مانند \lr{PRM}، \lr{RRT} و \lr{RRT*} از نمونه‌های شناخته‌شده این خانواده هستند%
\cite{Kavraki1996PRM,LaValle1998RRT,Karaman2011RRTstar}.
این الگوریتم‌ها به‌ویژه در فضاهای پیکربندی با ابعاد بالا یا محیط‌هایی که تولید شبکه منظم در آن‌ها دشوار و پرهزینه است، کاربرد زیادی دارند%
\cite{LaValle2006Planning,Choset2005Principles}.

مزیت اصلی روش‌های نمونه‌برداری، توانایی آن‌ها در یافتن مسیر در محیط‌های پیچیده بدون نیاز به گسسته‌سازی کامل فضای پیکربندی است. برای نمونه، روش \lr{RRT} با رشد سریع در فضای آزاد، امکان یافتن مسیرهای قابل اجرا را در زمان نسبتاً کوتاه فراهم می‌کند و نسخه‌هایی مانند \lr{RRT*} تلاش می‌کنند کیفیت مسیر را از نظر بهینگی بهبود دهند%
\cite{LaValle1998RRT,Karaman2011RRTstar}.
با این حال، مسیرهای تولیدشده توسط این روش‌ها معمولاً خام، ناصاف و نیازمند هموارسازی یا پردازش تکمیلی هستند. همچنین، اگرچه این روش‌ها برای برنامه‌ریزی مسیر مفیدند، اما در حالت پایه برای محیط‌های به‌شدت دینامیکی و تصمیم‌گیری بلادرنگ در حضور موانع متحرک کافی نیستند%
\cite{LaValle2006Planning,Choset2005Principles}.

از منظر پژوهش حاضر، روش‌های نمونه‌برداری می‌توانند برای تولید مسیر اولیه یا مسیر مرجع مفید باشند، اما به‌تنهایی مسئله ناوبری پیوسته و ایمن اسکوتر در محیط دینامیکی را حل نمی‌کنند. دلیل اصلی آن است که اسکوتر نه‌تنها باید از نظر هندسی به هدف برسد، بلکه باید فرمان‌هایی تولید کند که با محدودیت‌های حرکتی، شتاب، ترمز، پایداری و نرمی حرکت نیز سازگار باشند.

\subsection{میدان‌های پتانسیل مصنوعی و مسئله کمینه‌های محلی}

یکی از نخستین رویکردهای مطرح در ناوبری ربات‌های متحرک، روش «میدان‌های پتانسیل مصنوعی%
\LTRfootnote{\lr{Artificial Potential Fields -- APF}}»
است. در این روش، هدف به‌عنوان یک پتانسیل جاذب و موانع به‌عنوان پتانسیل‌های دافع مدل می‌شوند و ربات با حرکت در راستای کاهش تابع پتانسیل به سمت هدف هدایت می‌شود. سادگی محاسبات، پیاده‌سازی آسان و امکان اجرای بلادرنگ از مهم‌ترین مزایای این روش به‌شمار می‌روند%
\cite{khatib1986}.

با وجود این مزایا، یکی از مشکلات بنیادین میدان‌های پتانسیل، گیر افتادن ربات در نقاط «کمینه محلی%
\LTRfootnote{\lr{Local Minima}}»
است. در چنین نقاطی، اثر پتانسیل جاذب هدف و پتانسیل‌های دافع موانع می‌تواند به تعادل برسد و گرادیان مؤثر میدان پتانسیل صفر یا بسیار کوچک شود؛ در نتیجه، ربات ممکن است بدون رسیدن به هدف متوقف شود. این پدیده به‌ویژه در محیط‌های دارای چینش‌های پیچیده موانع تشدید می‌شود. افزون بر این، روش میدان پتانسیل در حالت پایه معمولاً قیود دینامیکی پلتفرم، محدودیت ترمز، محدودیت شتاب و رفتار زمانی موانع متحرک را به‌صورت صریح در نظر نمی‌گیرد%
\cite{Koren1991PotentialFields,LaValle2006Planning,Choset2005Principles}.
بنابراین، اگرچه این روش از نظر مفهومی ساده و قابل فهم است، اما برای ناوبری ایمن در محیط‌های دینامیکی و برای پلتفرم‌هایی با دینامیک جدی، به‌تنهایی کافی نیست.

\subsection{روش‌های اجتناب از برخورد مبتنی بر سرعت}

برای مواجهه با موانع متحرک، خانواده‌ای از روش‌ها توسعه یافته‌اند که به‌جای تمرکز صرف بر موقعیت هندسی، فضای سرعت‌های مجاز و غیرمجاز را تحلیل می‌کنند. روش‌هایی مانند «مانع سرعت%
\LTRfootnote{\lr{Velocity Obstacle -- VO}}»،
«مانع سرعت متقابل%
\LTRfootnote{\lr{Reciprocal Velocity Obstacle -- RVO}}»
و «اجتناب بهینه متقابل از برخورد%
\LTRfootnote{\lr{Optimal Reciprocal Collision Avoidance -- ORCA}}»
در این دسته قرار می‌گیرند%
\cite{Fiorini1998VO,VanDenBerg2008RVO,VanDenBerg2011ORCA}.
ایده اصلی این روش‌ها آن است که با توجه به سرعت نسبی ربات و موانع متحرک، می‌توان نواحی‌ای از فضای سرعت را مشخص کرد که انتخاب آن‌ها در یک افق زمانی محدود به برخورد منجر می‌شود. سپس فرمان حرکتی ربات به‌گونه‌ای انتخاب می‌شود که از این نواحی خطرناک در فضای سرعت دوری کند.

این دسته از روش‌ها برای محیط‌های چندعاملی و موانع متحرک اهمیت زیادی دارند، زیرا مسئله برخورد را به‌صورت زمانی بررسی می‌کنند و فقط به فاصله لحظه‌ای از مانع وابسته نیستند. با این حال، عملکرد آن‌ها معمولاً به فرض‌هایی مانند حرکت خطی کوتاه‌مدت عامل‌ها، واکنش متقابل سایر عامل‌ها و دقت در تخمین سرعت موانع وابسته است. در محیط‌های واقعی، به‌ویژه در حضور انسان‌ها، رفتار موانع لزوماً خطی، منطقی یا قابل پیش‌بینی نیست. همچنین، این روش‌ها معمولاً مسئله را در سطح اجتناب محلی از برخورد حل می‌کنند و برای تصمیم‌گیری بلندمدت، انتخاب مسیر کلی و لحاظ کامل قیود دینامیکی پلتفرم نیازمند ترکیب با سایر روش‌ها هستند%
\cite{Fiorini1998VO,VanDenBerg2008RVO,VanDenBerg2011ORCA,LaValle2006Planning}.

\subsection{روش پنجره دینامیکی و برنامه‌ریزی در فضای سرعت}

روش «پنجره دینامیکی%
\LTRfootnote{\lr{Dynamic Window Approach -- DWA}}»
یکی از شناخته‌شده‌ترین روش‌های برنامه‌ریزی محلی برای ربات‌های متحرک است. برخلاف روش‌های مبتنی بر موقعیت، \lr{DWA} مستقیماً در فضای سرعت خطی و زاویه‌ای ربات جست‌وجو می‌کند و تنها سرعت‌هایی را در نظر می‌گیرد که با قیود دینامیکی پلتفرم، از جمله حداکثر سرعت، شتاب و توان توقف، سازگار باشند. برای هر زوج سرعت، یک مسیر کوتاه‌مدت شبیه‌سازی شده و سپس با استفاده از تابع هزینه، مناسب‌ترین فرمان انتخاب می‌شود%
\cite{fox1997dwa}.

محبوبیت \lr{DWA} ناشی از امکان اجرای بلادرنگ، سادگی نسبی و قابلیت لحاظ برخی قیود حرکتی است. با این حال، عملکرد این روش به‌شدت به تنظیم ضرایب تابع هزینه وابسته است؛ ضرایبی که معمولاً وزن‌هایی برای حرکت به سمت هدف، حفظ فاصله ایمن از موانع و انتخاب سرعت مناسب تعیین می‌کنند. در برخی پیاده‌سازی‌ها، معیارهایی مانند همواری حرکت یا کاهش تغییرات ناگهانی فرمان نیز به تابع هزینه افزوده می‌شود. در صورت تغییر محیط، سرعت ربات یا ویژگی‌های دینامیکی پلتفرم، این ضرایب ممکن است نیازمند تنظیم مجدد باشند%
\cite{fox1997dwa,LaValle2006Planning,Choset2005Principles}.

افزون بر این، افق تصمیم‌گیری \lr{DWA} کوتاه‌مدت است؛ بنابراین این روش معمولاً تصمیم محلی مناسبی تولید می‌کند، اما به‌تنهایی تضمین‌کننده تصمیم‌گیری بلندمدت یا انتخاب مسیر کلی بهینه نیست. به همین دلیل، در محیط‌های شلوغ و دینامیکی ممکن است رفتارهایی مانند توقف‌های مکرر، حرکت محافظه‌کارانه یا نوسان میان چند فرمان نزدیک به هم رخ دهد. این مسئله برای اسکوتر چهار چرخ محرک مهم‌تر است، زیرا تغییرات ناگهانی فرمان یا توقف‌های مکرر می‌تواند موجب کاهش نرمی حرکت، افت راحتی سرنشین و افزایش ریسک ایمنی شود.

\subsection{برنامه‌ریزی مسیر محلی مبتنی بر بهینه‌سازی و روش \lr{TEB}}

در سامانه‌های رباتیک عملی، علاوه بر روش‌هایی مانند \lr{DWA}، از برنامه‌ریزهای محلی مبتنی بر بهینه‌سازی نیز استفاده می‌شود. یکی از نمونه‌های رایج این دسته، روش «نوار کشسان زمان‌بندی‌شده%
\LTRfootnote{\lr{Timed Elastic Band -- TEB}}»
است. این روش را می‌توان توسعه‌ای از ایده نوارهای کشسان در برنامه‌ریزی حرکت دانست که در آن مسیر محلی به‌صورت مجموعه‌ای از حالت‌های متوالی همراه با اطلاعات زمانی نمایش داده می‌شود%
\cite{Quinlan1993ElasticBands,Rosmann2013TEB,Rosmann2015TEB}.
سپس با در نظر گرفتن معیارهایی مانند فاصله از موانع، محدودیت سرعت، محدودیت شتاب، زمان حرکت و برخی قیود سینماتیکی، یک مسئله بهینه‌سازی محلی حل می‌گردد. مزیت اصلی این روش آن است که برخلاف مسیرهای صرفاً هندسی، زمان‌بندی حرکت و برخی محدودیت‌های حرکتی را نیز در فرآیند برنامه‌ریزی وارد می‌کند.

با وجود این مزایا، روش‌هایی مانند \lr{TEB} نیز معمولاً به تابع هزینه، وزن‌دهی مناسب قیود، کیفیت اطلاعات محیط و سرعت به‌روزرسانی نقشه محلی وابسته‌اند. در محیط‌های دینامیکی، عملکرد این روش به توانایی سامانه در تخمین یا پیش‌بینی حرکت موانع و بازبهینه‌سازی سریع مسیر بستگی دارد. اگر رفتار موانع غیرقابل پیش‌بینی باشد یا داده‌های حسگری دارای تأخیر و نویز باشند، کیفیت مسیر بهینه‌شده و قابلیت اطمینان تصمیم محلی کاهش می‌یابد. بنابراین، اگرچه \lr{TEB} نسبت به برنامه‌ریزی هندسی ساده پیشرفته‌تر است و برای برنامه‌ریزی محلی ربات‌های متحرک کاربرد زیادی دارد، اما در محیط‌های بسیار پویا، شلوغ و ناقصاً مشاهده‌پذیر همچنان با محدودیت مواجه می‌شود%
\cite{Rosmann2013TEB,Rosmann2015TEB,LaValle2006Planning}.

\subsection{کنترل پیش‌بینانه مبتنی بر مدل و برنامه‌ریزی بهینه}

«کنترل پیش‌بینانه مبتنی بر مدل%
\LTRfootnote{\lr{Model Predictive Control -- MPC}}»
به‌عنوان یک چارچوب قدرتمند برای لحاظ کردن قیود حرکتی، قیود کنترلی و اهداف بهینه‌سازی در سامانه‌های دینامیکی مطرح شده است. در روش \lr{MPC}، یک مدل از دینامیک سیستم در افق زمانی محدود در نظر گرفته می‌شود و با حل یک مسئله بهینه‌سازی مقید، توالی مناسبی از فرمان‌های کنترلی تولید می‌گردد. سپس تنها فرمان اول اعمال شده و در گام زمانی بعد، مسئله با استفاده از اطلاعات جدید دوباره حل می‌شود. این سازوکار بازخوردی باعث می‌شود \lr{MPC} بتواند در صورت دسترسی به مدل و اطلاعات به‌روز، نسبت به تغییرات وضعیت سیستم و محیط واکنش نشان دهد%
\cite{rawlings2009mpc,Bemporad1999}.

مزیت اصلی \lr{MPC} توانایی آن در لحاظ هم‌زمان قیودی مانند محدودیت سرعت، محدودیت شتاب، محدودیت ورودی‌های کنترلی، فاصله ایمن از موانع و اهداف حرکتی است. این ویژگی برای پلتفرم‌هایی مانند اسکوتر چهار چرخ محرک اهمیت زیادی دارد، زیرا این نوع پلتفرم نسبت به ربات‌های کوچک آزمایشگاهی دارای اینرسی بیشتر، محدودیت توقف جدی‌تر و حساسیت بالاتر نسبت به تغییرات ناگهانی فرمان است. با این حال، عملکرد \lr{MPC} به دقت مدل دینامیکی، کیفیت پیش‌بینی وضعیت محیط و تنظیم مناسب تابع هزینه وابسته است. علاوه بر این، حل مسئله بهینه‌سازی در هر گام زمانی می‌تواند هزینه محاسباتی قابل توجهی ایجاد کند، به‌ویژه زمانی که قیود غیرخطی، موانع متحرک یا افق پیش‌بینی بلندتر در نظر گرفته شوند.

نسخه‌های مقاوم‌تر \lr{MPC}، مانند \lr{Tube MPC}، می‌توانند تا حدی اثر عدم قطعیت‌ها و اغتشاش‌های محدود را در نظر بگیرند. با این حال، طراحی قیود مقاوم، انتخاب ناحیه‌های مجاز، تنظیم وزن‌های تابع هزینه و تضمین عملکرد مناسب در سناریوهای متنوع همچنان فرآیندی پیچیده و وابسته به مدل و دانش طراح باقی می‌ماند%
\cite{Mayne2005TubeMPC}.
از این رو، هرچند \lr{MPC} از نظر نظری و عملی یکی از قوی‌ترین روش‌های کلاسیک برای کنترل و برنامه‌ریزی مقید محسوب می‌شود، اما استفاده از آن به‌عنوان راهکار مستقل برای ناوبری کاملاً تطبیقی در محیط‌های دینامیکی، شلوغ و ناقصاً قابل پیش‌بینی با چالش‌های مهمی همراه است.

\subsection{جمع‌بندی محدودیت‌های روش‌های کلاسیک}

مرور روش‌های کلاسیک نشان می‌دهد که هر خانواده از الگوریتم‌ها بخشی از مسئله ناوبری را به‌خوبی پوشش می‌دهد، اما هیچ‌کدام به‌تنهایی پاسخ کامل و عمومی برای ناوبری در محیط‌های دینامیکی و پیچیده ارائه نمی‌کنند. روش‌های جست‌وجومحور و نمونه‌برداری برای تولید مسیر سراسری مفیدند، اما معمولاً دینامیک پلتفرم و رفتار زمانی موانع را به‌صورت کامل لحاظ نمی‌کنند%
\cite{s25113394,alvarez2024DRLsurvey}. %
 روش‌های میدان پتانسیل ساده و سریع هستند، اما با مشکل کمینه محلی و رفتارهای ناپایدار مواجه می‌شوند. روش‌های مبتنی بر سرعت مانند \lr{DWA}، \lr{Velocity Obstacle} و \lr{ORCA} برای اجتناب محلی از برخورد مناسب‌اند، اما افق تصمیم‌گیری محدود و وابستگی به فرضیات ساده درباره حرکت موانع دارند. روش‌های بهینه‌سازی‌محور مانند \lr{TEB} و \lr{MPC} نیز امکان لحاظ قیود حرکتی را فراهم می‌کنند، اما به مدل، تابع هزینه و تنظیم دقیق پارامترها وابسته‌اند.

این محدودیت‌ها در مسئله ناوبری اسکوتر چهار چرخ محرک برجسته‌تر می‌شوند، زیرا تصمیم‌های حرکتی باید علاوه بر اجتناب از برخورد، با قیود دینامیکی، اینرسی، محدودیت شتاب و ترمز، نرمی حرکت و ایمنی سرنشین نیز سازگار باشند. بنابراین، مسئله پژوهش حاضر صرفاً یافتن یک مسیر هندسی از مبدأ به هدف نیست، بلکه یادگیری یا تولید سیاستی است که بتواند در محیط دینامیکی، فرمان‌های پیوسته، ایمن و قابل اجرا تولید کند.

\section{روش‌های یادگیری ماشین در ناوبری ربات‌های متحرک}

در ناوبری ربات‌های متحرک، تصمیم حرکتی بر اساس داده‌های حسگری و وضعیت محیط تولید می‌شود. روش‌های کلاسیک این تصمیم را با مدل‌های تحلیلی، قوانین از پیش تعریف‌شده و توابع هزینه می‌سازند، اما مدل‌سازی دقیق محیط و پیش‌بینی همه موانع در سناریوهای دینامیکی و شلوغ دشوار یا ناممکن است. در چنین شرایطی، یادگیری ماشین امکان استخراج الگوهای تصمیم‌گیری از داده و تجربه را فراهم می‌کند.

در این رویکرد، سامانه به‌جای اتکا به قوانین صریح، رابطه میان ورودی‌هایی مانند داده‌های دوربین و حسگرهای اینرسی و خروجی‌هایی مانند مسیر یا فرمان حرکتی را می‌آموزد%
\cite{Zhu2021}. %
این ویژگی می‌تواند وابستگی به مدل دقیق و تنظیم دستی پارامترها را کاهش دهد، به‌ویژه هنگامی که محیط متغیر، نامطمئن و پیش‌بینی‌ناپذیر است.

بااین‌حال، یادگیری ماشین شامل خانواده‌هایی با اهداف و قابلیت‌های متفاوت است؛ برخی برای پیش‌بینی و استخراج ویژگی مناسب‌اند و برخی تصمیم‌گیری و بهینه‌سازی رفتار در طول زمان را ممکن می‌کنند. ازاین‌رو، انتخاب روش متناسب با مسئله ناوبری نقش تعیین‌کننده‌ای در عملکرد سامانه دارد.

\subsection{یادگیری نظارتی در ناوبری ربات‌های متحرک}

«یادگیری نظارتی%
\LTRfootnote{Supervised Learning}»
با استفاده از داده‌های برچسب‌خورده، نگاشتی میان ورودی و خروجی را با کمینه‌کردن خطای پیش‌بینی می‌آموزد. در ناوبری، ورودی‌ها معمولاً داده‌های حسگری و اطلاعات هدف و خروجی‌ها مسیر محلی یا فرمان‌های حرکتی هستند.

در رویکردهای انتها به انتها، تصمیم ناوبری مستقیماً از ادراک به کنترل نگاشت می‌شود. برای نمونه، «پفایفر%
\LTRfootnote{Pfeiffer}»%
و همکاران چارچوبی ارائه کردند که داده‌های لایدار و موقعیت هدف را به سرعت خطی و زاویه‌ای ربات تبدیل می‌کند. داده‌های آموزشی این مدل از اجرای یک سیاست مرجع مبتنی بر برنامه‌ریزی کلاسیک جمع‌آوری شد و شبکه عصبی رفتار آن را تقلید کرد%
\cite{pfeiffer2018nav}.

مزیت این روش، سادگی تصمیم‌گیری و سرعت اجرای بالا پس از آموزش است؛ زیرا به حل بهینه‌سازی برخط نیاز ندارد و می‌تواند از داده انسانی یا الگوریتم مرجع برای تولید رفتارهای هموار در محیط‌های نسبتاً ساخت‌یافته استفاده کند. بااین‌حال، عملکرد آن به کیفیت و پوشش داده‌های آموزشی وابسته است و شرایط دیده‌نشده ممکن است رفتار نامناسب یا ناایمن ایجاد کند. این مسئله با عنوان «انتقال توزیع%
\LTRfootnote{distribution shift}»% 
شناخته می‌شود و در محیط‌های دینامیکی اهمیت بیشتری دارد.

محدودیت دیگر، ماهیت غیرفعال یادگیری است؛ مدل تنها از داده‌های ثابت می‌آموزد و پیامد تصمیم‌های خود را تجربه نمی‌کند. بنابراین، سازوکار مستقیمی برای بهینه‌سازی هدف بلندمدت یا اصلاح رفتار از طریق تعامل ندارد و بیشتر برای تقلید رفتار یا پیش‌بینی محلی مناسب است تا تصمیم‌گیری تطبیقی در محیط پویا%
\cite{alvarez2024DRLsurvey}.
در مجموع، یادگیری نظارتی برای استخراج الگوهای حرکتی و مدل‌سازی رفتار مرجع مفید است، اما به‌تنهایی نیاز ناوبری خودران در شرایط غیرقابل پیش‌بینی را برآورده نمی‌کند.

\subsection{یادگیری بدون نظارت در ناوبری ربات‌های متحرک}

«یادگیری بدون نظارت%
\LTRfootnote{Unsupervised Learning}» %
بدون برچسب صریح، ساختارها و الگوهای نهفته داده را استخراج می‌کند. در ناوبری، این روش بیشتر برای ادراک محیط، کاهش پیچیدگی داده‌های حسگری و یادگیری نمایش فشرده از وضعیت ربات و محیط به‌کار می‌رود. داده‌های حسگری معمولاً پربعد، پرنویز و پیچیده‌اند و استخراج ویژگی‌های معنادار می‌تواند هزینه محاسباتی را کاهش و ورودی مناسب‌تری برای تصمیم‌گیری فراهم کند%
\cite{Zhu2021,s25113394}. 

برخی پژوهش‌ها این نقش را به تولید مستقیم مسیر گسترش داده‌اند. برای نمونه، در چارچوب %
 \lr{UP3}%
، مدل بدون برچسب دستی یا سیاست مرجع، با استفاده از ساختار داده‌های مشاهده‌شده و پیش‌بینی تحول محیط، مسیرهای قابل اجرا و ایمن را در شرایط مشاهده محدود تولید می‌کند. نتایج نشان داده‌اند که این روش می‌تواند بدون تعامل گسترده یا آموزش مبتنی بر پاداش، مسیرهایی هموار و سازگار با محدودیت‌های حرکتی ایجاد کند%
\cite{Luo2025UP3}.

این رویکرد هنگامی که داده برچسب‌خورده یا طراحی تابع پاداش دشوار است، می‌تواند مکمل یا جایگزین مناسبی باشد. بااین‌حال، معمولاً به فرض‌هایی درباره ساختار یا پیش‌بینی‌پذیری محیط متکی است و تعمیم‌پذیری، تضمین ایمنی و سازگاری آن با تغییرات سریع همچنان محدود است. ازاین‌رو، یادگیری بدون نظارت در بیشتر کاربردهای ناوبری نقش مکمل سایر روش‌های تصمیم‌گیری و کنترل را دارد.

\subsection{یادگیری تقلیدی در ناوبری ربات‌های متحرک}

«یادگیری تقلیدی%
\LTRfootnote{Imitation Learning}»
رفتار یک عامل مرجع، معمولاً انسان یا الگوریتم کنترلی، را از دنباله مشاهدات و اعمال متناظر می‌آموزد. رفتار مطلوب به‌جای تابع پاداش صریح، از نمایش‌های ارائه‌شده استخراج می‌شود. این روش از یک سو مانند یادگیری نظارتی بر داده برچسب‌خورده تکیه دارد و از سوی دیگر مستقیماً با انتخاب عمل سروکار دارد؛ ازاین‌رو برای رفتارهایی که طراحی صریح آن‌ها دشوار است، مانند اجتناب از برخورد در محیط‌های شلوغ، جذاب است%
\cite{argall2009survey}.

در پیاده‌سازی رایج، یک سیاست کلاسیک یا انسانی دنباله‌های حرکتی مرجع را تولید می‌کند و ربات نگاشت وضعیت محیط به فرمان را می‌آموزد. این رویکرد دانش روش‌های کلاسیک را با انعطاف مدل‌های یادگیری ترکیب می‌کند و می‌تواند رفتارهایی هموار در محیط‌های نسبتاً پیچیده ایجاد کند. در رویکرد انتها به انتها نیز «کادویلا%
 \LTRfootnote{Codevilla}»%
و همکاران با داده‌های رانندگی انسانی، رفتار ناوبری را مستقیماً از داده‌های ادراکی آموختند. هرچند این مطالعه در زمینه رانندگی خودکار انجام شد، چارچوب آن نشان می‌دهد که یادگیری تقلیدی می‌تواند رفتارهای پیچیده را بدون قوانین کنترلی صریح بازتولید کند%
\cite{codevilla2018endtoend}.

محدودیت اصلی، وابستگی سیاست به کیفیت و تنوع نمایش‌هاست. رفتار مرجع ناقص یا داده محدود، در شرایط جدید افت عملکرد و مشکل \lr{covariate shift} ایجاد می‌کند. افزون بر آن، تقلید صرف سازوکار صریحی برای بهینه‌سازی هدف بلندمدت، بهبود تدریجی عملکرد یا کشف راهبردهایی فراتر از داده مرجع ندارد. بنابراین، با وجود ارزش آن به‌عنوان نقطه آغاز یادگیری، به‌تنهایی برای رفتار تطبیق‌پذیر و بهینه در محیط پویا کافی نیست%
\cite{Zhu2021}.


\section{یادگیری تقویتی به‌عنوان چارچوب تصمیم‌گیری در ناوبری ربات‌های متحرک}

یادگیری تقویتی از طریق تعامل مستقیم عامل با محیط، رفتار را بر اساس پیامد اعمال اصلاح می‌کند و برخلاف یادگیری نظارتی یا تقلیدی، به خروجی‌های برچسب‌خورده یا بازتولید صرف رفتار مرجع وابسته نیست. این چارچوب برای مسائلی مناسب است که تصمیم‌ها ماهیتی زمانی دارند و عمل فعلی، وضعیت و امکان‌های تصمیم‌گیری آینده را تغییر می‌دهد؛ ویژگی‌ای که در ناوبری ربات‌های متحرک به‌طور مستقیم وجود دارد.

در این رویکرد، هدف‌های بلندمدت مانند رسیدن به مقصد، اجتناب از برخورد و همواری حرکت با معیارهای کلی عملکرد تعریف می‌شوند و نیازی به تعیین خروجی مطلوب برای تمام وضعیت‌های ممکن نیست%
\cite{Zhu2021}.
عامل همچنین با تجربه مستقیم محیط می‌تواند راهبردهایی را کشف کند که در داده‌های از پیش جمع‌آوری‌شده یا رفتار مرجع وجود ندارند و در برابر شرایط غیرمنتظره سازگارتر باشد.

به همین دلیل، یادگیری تقویتی به یکی از چارچوب‌های اصلی تصمیم‌گیری در ناوبری تبدیل شده است و امکان ترکیب چند هدف حرکتی در یک سیاست منسجم را فراهم می‌کند. بااین‌حال، استفاده مؤثر از آن به انتخاب الگوریتم مناسب، لحاظ قیود حرکتی و مدیریت چالش‌های آموزش و اجرا وابسته است%
\cite{s25113394}.
\subsection{خانواده‌های الگوریتم‌های یادگیری تقویتی در ناوبری ربات‌های متحرک}

الگوریتم‌های یادگیری تقویتی مورد استفاده در ناوبری ربات‌های متحرک را می‌توان بر اساس نحوه‌ی نمایش سیاست و سازوکار یادگیری به چند خانواده‌ی اصلی تقسیم کرد. این دسته‌بندی در ادبیات پژوهشی به‌عنوان یک چارچوب مفهومی برای تحلیل مزایا و محدودیت‌های روش‌های مختلف  یادگیری تقویتی در مسائل ناوبری شناخته می‌شود و نقش مهمی در انتخاب الگوریتم مناسب برای سامانه‌های حرکتی واقعی ایفا می‌کند. مرورهای جامع اخیر نشان می‌دهند که هر یک از این خانواده‌ها، رفتار متفاوتی در مواجهه با پیچیدگی‌های محیط‌های دینامیکی از خود نشان می‌دهند%
\cite{Alsolami2025AutonomousNavigationSurvey,Le2024DRLNavigationSurvey}.
\begin{itemize}
\item{\textbf{روش‌های مبتنی بر ارزش:}}
نخستین خانواده، روش‌های «\textbf{مبتنی بر ارزش}%
\LTRfootnote{Value-based}» هستند که در آن‌ها سیاست به‌طور غیرمستقیم و از طریق یادگیری تابع ارزش یا تابع $Q$ استخراج می‌شود. الگوریتم‌هایی مانند 
\lr{Q-learning} و نسخه‌های عمیق‌تر آن نظیر DQN در این دسته قرار می‌گیرند. این روش‌ها عمدتاً برای مسائل با فضای عمل گسسته طراحی شده‌اند و در برخی سناریوهای ساده‌ی ناوبری نتایج قابل قبولی ارائه داده‌اند. با این حال، گسسته‌سازی فضای عمل در سامانه‌های حرکتی واقعی می‌تواند منجر به کاهش دقت کنترلی و بروز رفتارهای ناپیوسته شود؛ مسئله‌ای که کاربرد این خانواده را در پلتفرم‌های دارای دینامیک پیوسته محدود می‌کند%
\cite{Alsolami2025AutonomousNavigationSurvey}.

\item{\textbf{روش‌های مبتنی بر گرادیان سیاست:}}
خانواده‌ی دوم شامل روش‌های «\textbf{مبتنی بر گرادیان سیاست}%
\LTRfootnote{Policy gradient}» است که سیاست کنترلی را به‌صورت مستقیم یاد می‌گیرند. الگوریتم‌هایی نظیر REINFORCE، TRPO و PPO در این دسته قرار دارند. این روش‌ها امکان کار با فضای عمل پیوسته را فراهم می‌کنند و از این نظر برای مسائل ناوبری ربات‌های متحرک مناسب‌تر از روش‌های مبتنی بر ارزش هستند. با این حال، بسیاری از این الگوریتم‌ها از نوع on-policy بوده و برای یادگیری مؤثر نیازمند داده‌های جدید و همگن هستند؛ موضوعی که می‌تواند به کاهش کارایی نمونه و افزایش هزینه‌ی آموزش در محیط‌های پیچیده و دینامیکی منجر شود%
\cite{Kober2013RLRobotics,Le2024DRLNavigationSurvey}.

\item{\textbf{روش‌های کنشگر--منتقد:}}
سومین خانواده، روش‌های کنشگر منتقد هستند که ترکیبی از دو رویکرد قبلی را به‌کار می‌گیرند. در این چارچوب، بخش کنشگر مسئول یادگیری سیاست و بخش منتقد مسئول برآورد کیفیت تصمیم‌ها از طریق تابع ارزش است. این ساختار امکان یادگیری پایدارتر در فضاهای حالت و عمل با ابعاد بالا را فراهم می‌کند و در سال‌های اخیر به‌طور گسترده در مسائل ناوبری ربات‌های متحرک مورد استفاده قرار گرفته است. الگوریتم‌های مختلفی در این خانواده معرفی شده‌اند که در پژوهش‌های جدید به‌عنوان گزینه‌های مناسب برای کنترل پیوسته و ناوبری در محیط‌های دینامیکی شناخته می‌شوند%
\cite{Alsolami2025AutonomousNavigationSurvey,Le2024DRLNavigationSurvey}.

یکی از ویژگی‌های مهم روش‌های کنشگر-منتقد، امکان استفاده از داده‌های off-policy و در نتیجه بهبود کارایی نمونه در مقایسه با بسیاری از روش‌های on-policy است. این ویژگی در مسائل ناوبری واقعی، که جمع‌آوری داده‌ی فیزیکی پرهزینه و همراه با ریسک است، اهمیت ویژه‌ای دارد. با این حال، این الگوریتم‌ها نیز با چالش‌هایی نظیر حساسیت به تنظیم هایپرپارامترها و پیچیدگی پیاده‌سازی همراه هستند که باید در طراحی سامانه‌های عملی مدنظر قرار گیرند.
\end{itemize}
به‌طور کلی، مرور ادبیات نشان می‌دهد که روند پژوهش‌ها در حوزه ناوبری ربات‌های متحرک، به‌تدریج از روش‌های مبتنی بر ارزش به سمت روش‌های کنشگر-منقد با قابلیت مدیریت فضای عمل پیوسته و قیود دینامیکی حرکت کرده است. این روند، زمینه‌ساز توسعه و به‌کارگیری چارچوب‌های پیشرفته‌ی یادگیری تقویتی در سامانه‌های حرکتی واقعی شده و مبنای تصمیم‌گیری در پژوهش‌های جدید را تشکیل می‌دهد%
\cite{Alsolami2025AutonomousNavigationSurvey}.
\subsection{الگوریتم‌های یادگیری تقویتی برای فضاهای عمل پیوسته در ناوبری}

در بسیاری از کاربردهای عملی ناوبری ربات‌های متحرک، فضای عمل به‌صورت ذاتی «پیوسته»%
\LTRfootnote{continuous action space}%
	است. برای مثال، فرمان‌های سرعت خطی و زاویه‌ای، گشتاور چرخ‌ها یا شتاب‌های اعمال‌شده به سامانه‌های حرکتی، همگی مقادیری پیوسته هستند. در چنین شرایطی، گسسته‌سازی فضای عمل، که در روش‌های مبتنی بر ارزش رایج است، می‌تواند منجر به کاهش دقت کنترلی، افزایش نوسانات حرکتی و رفتارهای غیرهموار شود. از این رو، بخش قابل توجهی از پژوهش‌های اخیر در ناوبری ربات‌ها به استفاده از الگوریتم‌های یادگیری تقویتی مناسب برای فضاهای عمل پیوسته معطوف شده است%
\cite{Kober2013RLRobotics,Alsolami2025AutonomousNavigationSurvey}.

یکی از نخستین تلاش‌های موفق در این زمینه، الگوریتم «گرادیان سیاست قطعی عمیق%
\LTRfootnote{Deep Deterministic Policy Gradient--DDPG}» است که به‌عنوان یک روش کنشگر-منتقد با سیاست قطعی معرفی شد%
\cite{Lillicrap2016}. %
 این الگوریتم امکان یادگیری مستقیم سیاست‌های پیوسته را فراهم می‌کند و به‌طور گسترده در مسائل کنترل ربات‌های متحرک مورد استفاده قرار گرفته است. با این حال، مطالعات نشان داده‌اند که DDPG نسبت به نویز، تنظیم هایپرپارامترها و برآورد ناپایدار تابع ارزش حساس است، که این مسئله می‌تواند در محیط‌های دینامیکی ناوبری منجر به رفتارهای ناپایدار شود%
\cite{Kober2013RLRobotics}.

برای رفع بخشی از این محدودیت‌ها، الگوریتم «گرادیان سیاست عمیق قطعی با تاخیر دوگانه%
\LTRfootnote{Twin Delayed Deep Deterministic Policy Gradient--TD3}» معرفی شد که با استفاده از دو شبکه‌ی منقد و به‌روزرسانی‌های تأخیری سیاست، مشکل بیش‌برآوردی تابع ارزش را کاهش می‌دهد%
\cite{FujimotoTD3}. %
 این بهبود باعث افزایش پایداری آموزش و کیفیت سیاست‌های یادگرفته‌شده در مقایسه با DDPG شده است و در برخی مطالعات ناوبری بدون نقشه عملکرد بهتری از خود نشان داده است. با این حال، سیاست‌های قطعی \lr{TD3} همچنان می‌توانند در محیط‌های بسیار پویا یا دارای عدم‌قطعیت بالا، از نظر اکتشاف فضای حالت با محدودیت مواجه شوند%
\cite{Nasti2025MaplessTD3}.

در مقابل، الگوریتم‌های مبتنی بر سیاست‌های «\textbf{تصادفی}%
\LTRfootnote{stochastic policies}» رویکرد متفاوتی را دنبال می‌کنند. در این میان، الگوریتم «\textbf{کنشگر-منتقد نرم}%
\LTRfootnote{Soft Actor--Critic-SAC}» با افزودن مؤلفه‌ی آنتروپی به تابع هدف، تلاش می‌کند میان بهره‌برداری از سیاست‌های آموخته‌شده و اکتشاف فضای حالت تعادل برقرار کند
\cite{Haarnoja2018SAC}. %
 این ویژگی به‌ویژه در محیط‌های ناوبری دینامیکی اهمیت دارد، زیرا ربات باید بتواند در مواجهه با موانع متحرک و تغییرات پیش‌بینی‌ناپذیر، رفتارهای متنوع و سازگار از خود نشان دهد. مطالعات مروری اخیر نشان می‌دهند که الگوریتم‌های مبتنی بر سیاست‌های تصادفی، از نظر پایداری آموزش و کارایی نمونه، مزایایی نسبت به سیاست‌های کاملاً قطعی دارند%
\cite{Le2024DRLNavigationSurvey,Alsolami2025AutonomousNavigationSurvey}.

به‌طور کلی، ادبیات پژوهشی نشان می‌دهد که حرکت از سیاست‌های گسسته به سیاست‌های پیوسته، و از روش‌های مبتنی بر ارزش به سمت چارچوب‌ها کنشگر--منقد پیشرفته، یک روند غالب در ناوبری ربات‌های متحرک بوده است. این روند به‌ویژه در سامانه‌هایی با قیود دینامیکی، نیاز به حرکت هموار و تعامل ایمن با محیط، اهمیت بیشتری پیدا می‌کند و مبنای استفاده از الگوریتم‌های پیشرفته‌ی یادگیری تقویتی در پژوهش‌های اخیر را شکل می‌دهد%
\cite{Alsolami2025AutonomousNavigationSurvey}.

\subsection{مزایای یادگیری تقویتی نسبت به روش‌های پیشین ناوبری}

روش‌های کلاسیک با وجود کارایی و تحلیل‌پذیری، معمولاً به مدل دقیق محیط یا فرض‌های ساده درباره رفتار موانع متکی‌اند و در شرایط متغیر یا با اطلاعات ناقص افت می‌کنند. یادگیری تقویتی با تعامل مستقیم، چارچوب انعطاف‌پذیرتری برای تصمیم‌گیری در چنین محیط‌هایی فراهم می‌کند%
\cite{Wang2020G2RL}.
برخلاف یادگیری نظارتی، تعیین خروجی بهینه برای همه وضعیت‌ها ضروری نیست و سیاست از معیارهای کلی مانند ایمنی، کارایی و رسیدن به هدف استخراج می‌شود؛ بنابراین عامل می‌تواند رفتاری فراتر از نگاشت‌های موجود در داده آموزشی بیاموزد%
\cite{Zhu2021}.

یادگیری تقلیدی نیز به داده‌های نمایشی وابسته است و در وضعیت‌های دیده‌نشده ممکن است افت کند. در مقابل، یادگیری تقویتی با تعامل فعال و اصلاح تدریجی سیاست، سازگاری بیشتری با تغییرات محیطی دارد؛ مزیتی که در محیط‌های دینامیکی و شلوغ اهمیت ویژه‌ای پیدا می‌کند%
\cite{s25113394}.

مزیت دیگر، در نظر گرفتن پیامدهای بلندمدت است. یک فرمان ممکن است در لحظه ایمن باشد، اما ربات را در گام‌های بعدی در وضعیت نامطلوب قرار دهد. بیشینه‌سازی عملکرد تجمعی، توالی تصمیم‌ها را یکپارچه ارزیابی می‌کند و وابستگی به واکنش‌های صرفاً محلی را کاهش می‌دهد%
\cite{Wang2020G2RL}.
همچنین اهدافی مانند اجتناب از برخورد، کوتاهی مسیر، همواری حرکت و محدودیت‌های حرکتی را می‌توان در یک فرآیند یادگیری ترکیب کرد و رفتار منسجم‌تری نسبت به مجموعه‌ای از قواعد جداگانه به دست آورد%
\cite{Lee2022DRLNavigation}.

در مجموع، انعطاف در برابر عدم قطعیت، توانایی سازگاری و لحاظ اهداف بلندمدت، یادگیری تقویتی را برای ناوبری پیچیده مناسب می‌کند؛ هرچند بهره‌گیری عملی از این مزایا مستلزم توجه به محدودیت‌های آموزش، ایمنی و تعمیم‌پذیری است.
\subsection{چالش‌ها و محدودیت‌های یادگیری تقویتی در ناوبری}

یادگیری تقویتی عمیق علی‌رغم پیشرفت‌های چشمگیر در زمینه ناوبری ربات‌های متحرک، هنوز با چالش‌های مهمی در کاربردهای واقعی مواجه است. برخلاف روش‌های کلاسیک که رفتار سیستم به‌صورت صریح و قابل پیش‌بینی طراحی می‌شود، در DRL سیاست کنترلی از طریق تجربه و تعامل با محیط شکل می‌گیرد؛ این ویژگی اگرچه انعطاف‌پذیری بالایی ایجاد می‌کند، اما هم‌زمان منجر به بروز محدودیت‌هایی در آموزش و پیاده‌سازی عملی می‌شود%
\cite{Le2024DRLNavigationSurvey}.

یکی از چالش‌های اساسی، کارایی نمونه است. الگوریتم‌های DRL معمولاً برای همگرا شدن به سیاست‌های قابل قبول نیازمند حجم بسیار زیادی از داده‌های تعاملی هستند. این مسئله در شبیه‌سازی قابل مدیریت است، اما در محیط‌های واقعی می‌تواند منجر به افزایش هزینه، استهلاک سخت‌افزار و حتی بروز رفتارهای ناایمن در حین یادگیری شود. در محیط‌های دینامیکی که شامل موانع متحرک و تعامل با انسان هستند، این چالش اهمیت بیشتری پیدا می‌کند، زیرا امکان آزمون و خطای گسترده عملاً وجود ندارد%
\cite{Le2024DRLNavigationSurvey}.

چالش مهم دیگر پایداری و همگرایی الگوریتم‌ها در فضاهای حالت و عمل با ابعاد بالا است. ناوبری ربات‌های متحرک معمولاً با فضای حالت پیوسته، داده‌های حسگری پرنویز و محدودیت‌های دینامیکی همراه است. در چنین شرایطی، الگوریتم‌های DRL ممکن است دچار نوسانات شدید در فرآیند آموزش شوند یا به سیاست‌هایی همگرا شوند که تنها در شرایط خاص عملکرد مناسبی دارند. حساسیت بالا به تنظیم هایپرپارامترها و ساختار شبکه عصبی، این مشکل را تشدید می‌کند%
\cite{Le2024DRLNavigationSurvey}.

مسئلهٔ طراحی تابع پاداش یکی از چالش‌برانگیزترین جنبه‌های یادگیری تقویتی در ناوبری است. در مسائل واقعی، ربات باید به‌صورت هم‌زمان چندین هدف متضاد را رعایت کند؛ برای مثال، حرکت سریع‌تر ممکن است با افزایش خطر برخورد همراه باشد، در حالی که حرکت ایمن‌تر می‌تواند منجر به مسیرهای طولانی و ناکارآمد شود. ترکیب این اهداف در قالب یک تابع پاداش یکتا بدون ایجاد سوگیری یا رفتارهای ناخواسته، نیازمند طراحی دقیق و آزمون‌های گسترده است. در بسیاری از مطالعات گزارش شده است که تعریف نامناسب پاداش می‌تواند باعث یادگیری رفتارهای غیرایمن یا غیرکاربردی شود%
\cite{Le2024DRLNavigationSurvey}.

چالش قابل توجه دیگر انتقال از شبیه‌سازی به واقعیت است. اگرچه شبیه‌سازی ابزار قدرتمندی برای آموزش الگوریتم‌های DRL فراهم می‌کند، اما تفاوت‌های اجتناب‌ناپذیر میان شبیه‌سازی و دنیای واقعی، از جمله نویز حسگرها، مدل‌سازی ناقص دینامیک و رفتار غیرقابل پیش‌بینی موانع، می‌تواند باعث کاهش شدید عملکرد سیاست‌ های آموزش‌دیده در محیط واقعی شود. این مسئله یکی از موانع اصلی در تجاری‌سازی و کاربرد صنعتی روش‌های DRL در ناوبری ربات‌ها محسوب می‌شود%
\cite{Le2024DRLNavigationSurvey}.

در سناریوهای «ناوبری بدون نقشه%
\LTRfootnote{mapless navigation}»، چالش‌های فوق تشدید می‌شوند. در این حالت، ربات بدون دسترسی به نقشهٔ از پیش‌ساخته‌شده باید تنها بر اساس داده‌های آنی حسگرها تصمیم‌گیری کند. این امر منجر به پاداش‌های پراکنده، دشواری در کاوش فضای حالت و افزایش احتمال همگرایی به سیاست‌های زیر‌بهینه می‌شود. مطالعات اخیر نشان داده‌اند که استفاده از الگوریتم‌های کشگر--منقد بهبود‌یافته، مانند نسخه‌های اصلاح‌شدهٔ \lr{TD3} می‌تواند تا حدی پایداری یادگیری و کیفیت سیاست‌ها را در چنین شرایطی افزایش دهد، هرچند چالش‌های بنیادین همچنان باقی می‌مانند%
\cite{Nasti2025MaplessTD3}.

در مجموع، چالش‌هایی مانند کارایی نمونه، پایداری آموزش، طراحی تابع پاداش و انتقال به محیط‌های واقعی نشان می‌دهند که یادگیری تقویتی عمیق، علی‌رغم توانمندی‌های بالقوه، هنوز نیازمند توسعه و تکامل بیشتر برای کاربرد مطمئن در ناوبری ربات‌های متحرک در محیط‌های دینامیکی است. این چالش‌ها زمینه‌ساز پژوهش‌های گسترده‌ای در سال‌های اخیر بوده‌اند و انگیزهٔ اصلی برای توسعه چارچوب‌های ترکیبی و الگوریتم‌های پیشرفته‌تر محسوب می‌شوند.
\subsection{رویکردهای ایمن و ترکیبی در ناوبری مبتنی بر یادگیری تقویتی}

با وجود توانمندی‌های یادگیری تقویتی عمیق در تولید سیاست‌های کنترلی تطبیقی، یکی از مهم‌ترین دغدغه‌های مطرح در ادبیات ناوبری ربات‌های متحرک، مسئله‌ی \textbf{ایمنی} و رعایت قیود فیزیکی در طول فرآیند یادگیری و اجرا است. برخلاف روش‌های کلاسیک کنترل که قیود سیستم به‌صورت صریح در طراحی لحاظ می‌شوند، الگوریتم‌های یادگیری تقویتی ذاتاً فاقد تضمین‌های ایمنی هستند و در صورت عدم اعمال ملاحظات اضافی، ممکن است به رفتارهای ناایمن یا غیرقابل پیش‌بینی منجر شوند%
\cite{Le2024DRLNavigationSurvey}.

\begin{itemize}
\item{\textbf{یادگیری تقویتی ایمن و مقید:}}
یکی از رویکردهای مطرح در ادبیات، «یادگیری تقویتی ایمن%
\LTRfootnote{Safe Reinforcement Learning}» یا «یادگیری تقویتی مقید%
%
\LTRfootnote{Constrained Reinforcement Learning}» است که در آن، قیود ایمنی به‌صورت صریح در فرآیند یادگیری وارد می‌شوند. این قیود می‌توانند شامل محدودیت‌هایی بر فاصله از موانع، سرعت و شتاب مجاز، یا سایر پارامترهای مرتبط با ایمنی حرکت باشند. در این چارچوب‌ها، هدف تنها بیشینه‌سازی پاداش نیست، بلکه رعایت قیود نیز به‌عنوان بخشی از مسئله‌ی تصمیم‌گیری در نظر گرفته می‌شود%
\cite{Alsolami2025AutonomousNavigationSurvey}.

\item{\textbf{طراحی پاداش با در نظر گرفتن ایمنی و کیفیت حرکت:}}
رویکرد رایج‌تر در بسیاری از کاربردهای ناوبری، استفاده از «شکل‌دهی پاداش%
\LTRfootnote{reward shaping}» برای اعمال ملاحظات ایمنی و کیفی است. در این روش، به‌جای تعریف قیود صریح، مؤلفه‌های تنبیهی در تابع پاداش گنجانده می‌شوند تا رفتارهایی مانند نزدیک شدن بیش از حد به موانع، تغییرات ناگهانی فرمان یا نوسانات شدید حرکتی کاهش یابند. این نوع طراحی پاداش به‌ویژه برای سامانه‌های حرکتی واقعی که نرمی حرکت، راحتی و پایداری اهمیت دارند، مورد توجه قرار گرفته است%
\cite{Le2024DRLNavigationSurvey}.

\item{\textbf{روش‌های ترکیبی مبتنی بر یادگیری تقویتی:}}
دسته‌ی دیگری از پژوهش‌ها به سمت «روش‌های ترکیبی%
\LTRfootnote{hybrid approaches}» حرکت کرده‌اند که در آن‌ها یادگیری تقویتی با روش‌های کلاسیک برنامه‌ریزی مسیر یا کنترل ترکیب می‌شود. در این رویکردها، یادگیری تقویتی معمولاً نقش تصمیم‌گیری سطح بالا یا تولید سیاست کلی را بر عهده دارد، در حالی که تضمین ایمنی، رعایت قیود دینامیکی و اجرای دقیق فرمان‌ها به ماژول‌های کلاسیک مانند برنامه‌ریزها یا کنترل‌کننده‌های پیش‌بین سپرده می‌شود. این ترکیب امکان بهره‌گیری از انعطاف‌پذیری RL و در عین حال حفظ پایداری و ایمنی روش‌های کلاسیک را فراهم می‌کند%
\cite{Alsolami2025AutonomousNavigationSurvey}.
\end{itemize}
به‌طور کلی، مرور ادبیات نشان می‌دهد که برای کاربردهای واقعی ناوبری، به‌ویژه در پلتفرم‌هایی با جرم، ابعاد و محدودیت‌های دینامیکی قابل توجه، استفاده از یادگیری تقویتی به‌صورت خام کمتر مورد استقبال قرار گرفته و اغلب با ملاحظات ایمنی، طراحی پاداش هدفمند یا چارچوب‌های ترکیبی همراه بوده است. این رویکردها بیانگر تلاش پژوهشگران برای کاهش فاصله‌ی میان قابلیت‌های بالقوه‌ی یادگیری تقویتی و الزامات عملی سامانه‌های حرکتی واقعی هستند.
\subsection{مقایسه مطالعات منتخب در حوزه ناوبری ربات‌های متحرک}
به‌منظور جمع‌بندی منسجم مطالعات مرورشده، جدول \ref{tab:nav_literature_comparison} مقایسه‌ای میان مهم‌ترین خانواده‌های روش‌های کلاسیک، روش‌های مبتنی بر یادگیری و مطالعات نزدیک به مسئله پژوهش حاضر ارائه می‌دهد. هدف از این جدول، صرفاً فهرست کردن منابع نیست، بلکه روشن کردن نقش هر گروه از روش‌ها، نقاط قوت آن‌ها و محدودیت‌های آن‌ها نسبت به مسئله ناوبری اسکوتر چهار چرخ محرک در محیط دینامیکی است.

\begingroup
\scriptsize
\setlength{\tabcolsep}{3pt}
\renewcommand{\arraystretch}{1.35}
\newcolumntype{P}[1]{>{\raggedright\arraybackslash}p{#1}}
\newcolumntype{C}[1]{>{\centering\arraybackslash}p{#1}}
\newcolumntype{R}[1]{>{\raggedleft\arraybackslash}p{#1}}

\begin{longtable}{|C{2.35cm}|P{2.35cm}|P{2.65cm}|P{3.45cm}|P{3.95cm}|}
	\caption{مقایسه رویکردها و مطالعات منتخب در حوزه ناوبری ربات‌های متحرک}
	\label{tab:nav_literature_comparison}\\
	\hline
	\centering\textbf{مرجع} &
	\centering\textbf{رویکرد یا الگوریتم} &
	\centering\textbf{نوع مسئله یا محیط} &
	\centering\textbf{نقطه قوت و نقش در ادبیات} &
	\centering\arraybackslash\textbf{محدودیت نسبت به پژوهش حاضر} \\
	\hline
	\endfirsthead
	
	\hline
	\textbf{مرجع} &
	\textbf{رویکرد یا الگوریتم} &
	\textbf{نوع مسئله یا محیط} &
	\textbf{نقطه قوت و نقش در ادبیات} &
	\textbf{محدودیت نسبت به پژوهش حاضر} \\
	\hline
	\endhead
	
	\hline
	\endfoot
	
	\cite{Dijkstra1959,Hart1968AStar,koenig2005dstar} &
	روش‌های جست‌وجومحور مانند \lr{Dijkstra}، \lr{A*} و \lr{D* Lite} &
	برنامه‌ریزی مسیر سراسری در محیط‌های دارای نقشه &
	تولید مسیر کم‌هزینه در نمایش گرافی محیط و مناسب برای محیط‌های ایستا یا نیمه‌ایستا &
	مسیر تولیدشده عمدتاً هندسی است و به‌تنهایی قیود دینامیکی، محدودیت ترمز، موانع متحرک و تصمیم‌گیری لحظه‌ای را پوشش نمی‌دهد. \\
	\hline
	
	\cite{Kavraki1996PRM,LaValle1998RRT,Karaman2011RRTstar} &
	روش‌های نمونه‌برداری مانند \lr{PRM}، \lr{RRT} و \lr{RRT*} &
	برنامه‌ریزی در فضای پیکربندی، به‌ویژه در محیط‌های پیچیده &
	توانایی یافتن مسیر در فضاهای پیکربندی پیچیده بدون نیاز به گسسته‌سازی کامل محیط &
	مسیرها معمولاً نیازمند هموارسازی و پردازش تکمیلی هستند و برای تصمیم‌گیری بلادرنگ در محیط‌های شدیداً دینامیکی کافی نیستند. \\
	\hline
	
	\cite{khatib1986,Koren1991PotentialFields} &
	میدان پتانسیل مصنوعی \lr{APF} &
	اجتناب محلی از مانع و حرکت به سمت هدف &
	سادگی محاسباتی، پیاده‌سازی آسان و قابلیت اجرای بلادرنگ &
	دارای مشکل کمینه محلی، حساس به چینش موانع و فاقد لحاظ صریح قیود دینامیکی و رفتار زمانی موانع متحرک است. \\
	\hline
	
	\cite{Fiorini1998VO,VanDenBerg2008RVO,VanDenBerg2011ORCA} &
	روش‌های مبتنی بر سرعت مانند \lr{VO}، \lr{RVO} و \lr{ORCA} &
	اجتناب از برخورد در حضور موانع متحرک و عامل‌های چندگانه &
	تحلیل برخورد در فضای سرعت و توجه به زمان برخورد، نه فقط فاصله لحظه‌ای &
	وابسته به فرض حرکت کوتاه‌مدت قابل پیش‌بینی، تخمین دقیق سرعت موانع و معمولاً محدود به اجتناب محلی از برخورد است. \\
	\hline
	
	\cite{fox1997dwa} &
	روش پنجره دینامیکی \lr{DWA} &
	برنامه‌ریزی محلی در فضای سرعت &
	لحاظ سرعت، شتاب و توان توقف در انتخاب فرمان و امکان اجرای بلادرنگ &
	افق تصمیم‌گیری کوتاه دارد و عملکرد آن به تنظیم تابع هزینه وابسته است؛ در محیط‌های شلوغ ممکن است رفتار محافظه‌کارانه یا نوسانی ایجاد کند. \\
	\hline
	
	\cite{Quinlan1993ElasticBands,Rosmann2013TEB,Rosmann2015TEB} &
	نوار کشسان زمان‌بندی‌شده \lr{TEB} &
	برنامه‌ریزی محلی مبتنی بر بهینه‌سازی &
	ورود زمان‌بندی حرکت، محدودیت سرعت، شتاب و فاصله از موانع در بهینه‌سازی مسیر محلی &
	به وزن‌دهی تابع هزینه، کیفیت نقشه محلی، داده‌های حسگری و توان بازبهینه‌سازی سریع در محیط‌های دینامیکی وابسته است. \\
	\hline
	
	\cite{rawlings2009mpc,Bemporad1999,Mayne2005TubeMPC} &
	کنترل پیش‌بینانه مبتنی بر مدل \lr{MPC} و \lr{Tube MPC} &
	کنترل و برنامه‌ریزی مقید در سامانه‌های دینامیکی &
	توانایی لحاظ هم‌زمان قیود حرکتی، کنترلی، ایمنی و اهداف بهینه‌سازی در افق پیش‌بینی &
	نیازمند مدل نسبتاً دقیق، پیش‌بینی مناسب محیط و حل مسئله بهینه‌سازی در هر گام زمانی است؛ در محیط‌های شلوغ و نامطمئن هزینه طراحی و محاسبات بالا می‌رود. \\
	\hline
	
	\cite{s23249672} &
	مقایسه روش‌های سنتی و مبتنی بر یادگیری تقویتی عمیق &
	ناوبری داخلی در سناریوهای دینامیکی &
	نشان دادن تفاوت عملکرد روش‌های کلاسیک و \lr{DRL} در شرایط دینامیکی و وابستگی انتخاب روش به ویژگی محیط &
	بیشتر نقش مقایسه‌ای دارد و چارچوب اختصاصی برای پلتفرم‌هایی با قیود حرکتی جدی مانند اسکوتر ارائه نمی‌کند. \\
	\hline
	
	\cite{pfeiffer2018nav} &
	برنامه‌ریزی انتهابه‌انتها مبتنی بر داده &
	ناوبری ربات زمینی از ادراک تا تصمیم &
	کاهش وابستگی به طراحی دستی ماژول‌های جداگانه و اتصال مستقیم ادراک به تصمیم حرکتی &
	وابستگی بالا به داده آموزشی و تعمیم‌پذیری محدود در سناریوهای خارج از توزیع؛ قیود دینامیکی پلتفرم به‌صورت مرکزی بررسی نمی‌شود. \\
	\hline
	
	\cite{argall2009survey,codevilla2018endtoend} &
	یادگیری تقلیدی و راندن شرطی انتهابه‌انتها &
	یادگیری سیاست حرکتی از داده یا نمایش خبره &
	مناسب برای استفاده از تجربه انسانی یا سیاست مرجع و کاهش نیاز به طراحی صریح قوانین تصمیم‌گیری &
	کیفیت سیاست به داده خبره وابسته است و در مواجهه با شرایط دیده‌نشده ممکن است دچار افت عملکرد شود؛ بهینه‌سازی بلندمدت پاداش مانند \lr{RL} در مرکز روش نیست. \\
	\hline
	
	\cite{Luo2025UP3} &
	برنامه‌ریزی پیش‌بین بدون‌نظارت \lr{UP3} &
	محیط‌های ناقصاً مشاهده‌پذیر &
	توجه به پیش‌بینی آینده و تصمیم‌گیری در شرایطی که مشاهده محیط کامل نیست &
	مستقیماً چارچوب کنترل پیوسته مبتنی بر پاداش برای پلتفرم دارای قیود دینامیکی ارائه نمی‌کند. \\
	\hline
	
	\cite{Wang2020G2RL} &
	یادگیری تقویتی با راهنمایی سراسری &
	برنامه‌ریزی مسیر در محیط‌های دینامیکی &
	ترکیب اطلاعات برنامه‌ریزی سراسری با تصمیم‌گیری یادگیری‌محور برای بهبود ناوبری &
	هرچند برای محیط‌های دینامیکی مفید است، تمرکز اصلی آن بر پلتفرم‌هایی مانند اسکوتر چهار چرخ با قیود نرمی حرکت، توقف و راحتی سرنشین نیست. \\
	\hline
	
	\cite{Nasti2025MaplessTD3} &
	ناوبری بدون نقشه با الگوریتم بهبودیافته \lr{TD3} &
	ناوبری \lr{Mapless} در فضای عمل پیوسته &
	نزدیک به مسئله پژوهش حاضر از نظر استفاده از یادگیری تقویتی عمیق و فضای عمل پیوسته &
	تمرکز اصلی بر ناوبری بدون نقشه است و ابعاد خاص پلتفرم اسکوتر، ایمنی سرنشین، نرمی فرمان و قیود حرکتی سنگین‌تر به‌صورت محوری بررسی نمی‌شود. \\
	\hline
	
	\cite{Pacini2024RLwheelchair,Cui2025OutdoorWheelchair} &
	ناوبری و کمک‌رانش برای پلتفرم‌های کمک‌حرکتی مانند ویلچر هوشمند &
	پلتفرم‌های انسانی/کمک‌حرکتی در محیط واقعی یا نیمه‌واقعی &
	از نظر ماهیت پلتفرم و اهمیت ایمنی کاربر به مسئله اسکوتر خودران نزدیک‌تر از ربات‌های کوچک آزمایشگاهی است &
	معمولاً بر معماری کمک‌رانندگی، نقشه هزینه یا اجتناب از مانع تمرکز دارد و چارچوب کامل یادگیری سیاست پیوسته برای اسکوتر چهار چرخ محرک را پوشش نمی‌دهد. \\
	\hline
	
\end{longtable}

\endgroup

بررسی جدول \ref{tab:nav_literature_comparison} نشان می‌دهد که روش‌های کلاسیک هر یک بخشی از مسئله را پوشش می‌دهند: روش‌های جست‌وجویی و نمونه‌برداری برای مسیر سراسری، \lr{VO}، \lr{ORCA} و \lr{DWA} برای اجتناب محلی، و \lr{TEB} و \lr{MPC} برای لحاظ برخی قیود حرکتی و زمانی مفیدند. بااین‌حال، وابستگی به مدل صریح، تابع هزینه، نقشه محلی، پیش‌بینی حرکت موانع یا مدل دینامیکی، پاسخ‌گویی مستقل آن‌ها را در محیط‌های دینامیکی، شلوغ و ناقصاً مشاهده‌پذیر محدود می‌کند.

روش‌های داده‌محور و یادگیری تقویتی عمیق با یادگیری سیاست از داده یا تعامل، بخشی از این محدودیت‌ها را کاهش می‌دهند. بااین‌حال، بسیاری از مطالعات بر ربات‌های ساده‌تر، محیط‌های شبیه‌سازی‌شده یا فرضیات ساده‌شده متمرکزند و قیود خاص اسکوتر چهار چرخ محرک، از جمله اینرسی، محدودیت توقف، نرمی فرمان و ایمنی سرنشین را به‌صورت مرکزی بررسی نمی‌کنند. بنابراین، هنوز به چارچوبی نیاز است که در محیط دینامیکی و فضای عمل پیوسته، رسیدن به هدف، اجتناب از برخورد، قیود حرکتی، فرمان نرم و ایمنی را هم‌زمان پوشش دهد.

\section{جمع‌بندی مرور ادبیات و تبیین شکاف پژوهشی}

مرور ادبیات و مقایسه جدول \ref{tab:nav_literature_comparison} نشان داد که روش‌های کلاسیک با وجود بلوغ نظری، تحلیل‌پذیری و پیاده‌سازی روشن، به مدل محیط، تابع هزینه، کیفیت نقشه یا پیش‌بینی حرکت موانع وابسته‌اند. این وابستگی، عملکرد آن‌ها را در محیط‌های شلوغ، ناقصاً مشاهده‌پذیر و دارای موانع متحرک، به‌ویژه برای پلتفرم‌های دارای قیود حرکتی جدی، محدود می‌کند.

در مقابل، یادگیری تقویتی عمیق امکان یادگیری سیاست‌های غیرخطی و تطبیقی را از طریق تعامل فراهم می‌کند و برای ناوبری بدون نقشه، اجتناب از برخورد، محیط‌های دینامیکی و فضای عمل پیوسته به‌کار رفته است. بااین‌حال، کارایی نمونه، پایداری آموزش، طراحی تابع پاداش، ایمنی رفتار و تعمیم به شرایط دیده‌نشده همچنان از چالش‌های اصلی کاربرد واقعی آن هستند.

بخش قابل توجهی از مطالعات موجود نیز بر ربات‌های سبک، محیط‌های ساده‌شده یا شبیه‌سازی تمرکز دارد و معمولاً رسیدن به هدف و اجتناب از برخورد را ارزیابی می‌کند، درحالی‌که نرمی فرمان، شتاب و ترمز، اینرسی، پایداری و ایمنی یا راحتی سرنشین کمتر به‌طور هم‌زمان بررسی شده‌اند. برای اسکوتر چهار چرخ محرک، تصمیم ناوبری باید علاوه بر انتخاب مسیر و جلوگیری از برخورد، به فرمان‌های پیوسته، نرم، ایمن و قابل اجرا منجر شود.

ازاین‌رو، شکاف اصلی در کمبود چارچوبی است که ناوبری پلتفرمی با قیود اجرایی جدی را در محیط دینامیکی و فضای عمل پیوسته بررسی و میان رسیدن به هدف، اجتناب از برخورد، محدودیت‌های حرکتی، کاهش تغییرات ناگهانی فرمان و ایمنی تعادل برقرار کند. مسئله این پژوهش بنابراین طراحی و ارزیابی سیاست تصمیم‌گیری پیوسته برای ناوبری ایمن و هموار اسکوتر خودران است.

بر این اساس، شکاف‌های پژوهشی اصلی که این پایان‌نامه بر آن‌ها متمرکز است، به‌صورت زیر خلاصه می‌شوند:
\begin{itemize}
	\item کمبود مطالعات مسئله‌محور درباره ناوبری خودران پلتفرم‌هایی مانند اسکوتر چهار چرخ محرک که دارای اینرسی، محدودیت توقف و قیود حرکتی جدی هستند؛
	\item تمرکز محدود بسیاری از پژوهش‌های موجود بر ارزیابی هم‌زمان موفقیت در رسیدن به هدف، اجتناب از برخورد، نرمی حرکت و قابلیت اجرای فرمان‌ها؛
	\item وابستگی بسیاری از روش‌های کلاسیک به مدل‌سازی صریح محیط، تنظیم تابع هزینه یا پیش‌بینی حرکت موانع در محیط‌های دینامیکی؛
	\item چالش باقی‌مانده در طراحی سیاست‌های یادگیری تقویتی برای فضای عمل پیوسته که بتوانند میان کارایی ناوبری، ایمنی و نرمی فرمان تعادل برقرار کنند.
\end{itemize}

با توجه به این شکاف‌ها، پژوهش حاضر چارچوب ناوبری مبتنی بر یادگیری تقویتی را برای اسکوتر چهار چرخ محرک در محیط دینامیکی و تحت قیود عملی حرکت توسعه و ارزیابی می‌کند. این چارچوب علاوه بر رسیدن به هدف و اجتناب از برخورد، کیفیت فرمان‌های تولیدشده و سازگاری آن‌ها با محدودیت‌های پلتفرم را در نظر می‌گیرد. فصل بعد، فرمول‌بندی مسئله، مدل‌سازی محیط، تعریف حالت و عمل، طراحی پاداش و روش پیشنهادی پژوهش را تشریح می‌کند.
